\section{Introducción}

En banca, las APIs conectan canales propios, \textit{fintechs} y socios (Bhatia, 2022). Para
describirlas de forma uniforme se usa \textbf{OpenAPI}: un archivo YAML o JSON con \textit{paths},
operaciones HTTP, \textit{schemas} y restricciones legibles por máquina (Casas et al., 2021; Santos
et al., 2020). Ese contrato revisan desarrollo, integradores y gobernanza antes de exponer un
servicio.

En paralelo, \textbf{BIAN} (\textit{Banking Industry Architecture Network}) organiza el negocio en
\textbf{\ServiceDomain{}s}: capacidades como pagos o cuentas, con \textit{behaviors},
\textit{business objects} y relaciones de control (Farzi, 2021). BIAN no sustituye el código de la
API; funciona como \textbf{referencia} para contrastar si la interfaz expuesta cuenta la misma
historia de negocio.

Llamamos \textbf{alineación} al grado en que un contrato OpenAPI bancario y un \ServiceDomain{} BIAN
emparejado coinciden en estructura y significado. El producto del estudio es un \textbf{modelo de
puntuación}: \SemanticScore{} (S), \StructuralScore{} (E), \CoverageScore{} (C) y
\AlignmentScore{} $= \alpha S + \beta E + \gamma C$ con $\alpha + \beta + \gamma = 1$ (detalle en
§1.4 y Fig.~\ref{fig:alignment-score}).

A nivel de sector, la problemática (§\ref{sec:problematica}) muestra síntomas observables:
fragmentación semántica, costos de integración y casos documentados como TSB (2018) o \textit{open
banking} británico (CMA, 2016). El diagnóstico subyacente (§\ref{sec:objeto-estudio}) abstrae dos
jerarquías comparables ---contrato OpenAPI frente a \ServiceDomain{} BIAN---. El problema
tecnológico (§\ref{sec:planteamiento}) exige un procedimiento reproducible para cuantificar esa
correspondencia antes del despliegue.

La estrategia del plan es \textbf{diseñar} ese procedimiento y el modelo S/E/C, documentados con
trazabilidad en la matriz de consistencia (Anexo~C). La construcción del artefacto y la validación
empírica quedan fuera de alcance (§\ref{sec:alcances}).

La \Cref{fig:contexto} resume el flujo: sector API bancario, referencia BIAN, contrato OpenAPI como
entrada y scores como salida.

\begin{figure}[H]
  \centering
  \includegraphics[width=0.92\textwidth]{Figura_01_Contexto}
  \caption{Contexto: arquitectura API bancaria, BIAN, OpenAPI y modelo S/E/C}
  \label{fig:contexto}
\end{figure}

En la \Cref{sec:problematica} se describe la situación real; en §\ref{sec:planteamiento} y
§\ref{sec:objeto-estudio}, el problema subyacente y el objeto de estudio. Los
capítulos~\ref{cap:objetivos}--\ref{cap:metodologia} desarrollan objetivos, marco teórico y diseño
metodológico; el Capítulo~5 administra el plan.
